% GENERATED FILE. Edit canonical lessons, then run npm run generate:books.
% canonical-source-hash: fnv1a64:01116c0d1b9bdf5a
% canonical-lessons: ML-C555-tarattu, ML-C555-kuttukutumbam, ML-C555-anukutumbam, ML-C555-verpiriyuka, ML-C555-inannuka

\chapter{Generations Under One Roof}
\label{ch:ml-generations-under-one-roof}
\hlchaptermodality{555}

\begin{chapteropening}
\textbf{By the end of this chapter:} \emph{I can say a lullaby, a joint family, a nuclear family, to split up and to be reconciled.}

Complete the last lesson of chapter 555: I can say a lullaby, a joint family, a nuclear family, to split up and to be reconciled.
\end{chapteropening}

\section[\texorpdfstring{\ml{താരാട്ട്} (tārāṭṭŭ)}{താരാട്ട് (tārāṭṭŭ)}]{\ml{താരാട്ട്} (tārāṭṭŭ) — a lullaby}

\label{lesson:ML-C555-tarattu}



\begin{quote}
Before the new one: say the Malayalam for a sister-in-law, then the Malayalam for a family member.
\end{quote}

\subsection*{You'll want to know: \ml{താരാട്ട്}}
\textbf{\ml{താരാട്ട്}} — \emph{tārāṭṭŭ} — \textquotedblleft{}a lullaby\textquotedblright{}.

\begin{quote}
\ml{മുത്തശ്ശി} \ml{പാടിയ} \ml{താരാട്ട്} \ml{കേട്ട്} \ml{കുഞ്ഞ്} \ml{ഉറങ്ങി}.
\end{quote}

\begin{quote}
\emph{muttaśśi pāṭiya tārāṭṭŭ kēṭṭŭ kuññŭ uṟaṅṅi.}
\end{quote}

\begin{quote}
The baby fell asleep to the lullaby grandmother sang.
\end{quote}

\ml{താരാട്ട്} \ml{പാടുക}, to sing a lullaby.

\begin{cousinweb}
\textbf{\ml{താരാട്ട്}} (\emph{tārāṭṭŭ}) looks native, but DEDR has no entry that clearly holds it, so its origin is not settled.
\end{cousinweb}

\begin{grammarlens}[title={What you've built}]
One more word for describing what you lived and felt.
\end{grammarlens}

\subsection*{Your turn}
\begin{itemize}
\raggedright
  \item \emph{Say it:} \emph{tārāṭṭŭ}
  \item \emph{Say it:} \emph{tārāṭṭŭ}, once more
  \item \emph{Say it:} \emph{muttaśśi pāṭiya tārāṭṭŭ kēṭṭŭ kuññŭ uṟaṅṅi.} — The baby fell asleep to the lullaby grandmother sang.
  \item \emph{Recall:} say the Malayalam for a sister-in-law, then the Malayalam for a family member, then say \emph{tārāṭṭŭ} again
\end{itemize}

\begin{tcolorbox}[breakable,colback=teal!4,colframe=teal!35!black,title={Before you move on}]
What does \ml{താരാട്ട്} mean? (\textquotedblleft{}A lullaby\textquotedblright{}.) Say it once more.
\end{tcolorbox}
\section[\texorpdfstring{\ml{കൂട്ടുകുടുംബം} (kūṭṭukuṭumbaṁ)}{കൂട്ടുകുടുംബം (kūṭṭukuṭumbaṁ)}]{\ml{കൂട്ടുകുടുംബം} (kūṭṭukuṭumbaṁ) — a joint family}

\label{lesson:ML-C555-kuttukutumbam}



\begin{quote}
Before the new one: say the Malayalam for a family member, then the Malayalam for a lullaby.
\end{quote}

\subsection*{You'll want to know: \ml{കൂട്ടുകുടുംബം}}
\textbf{\ml{കൂട്ടുകുടുംബം}} — \emph{kūṭṭukuṭumbaṁ} — \textquotedblleft{}a joint family\textquotedblright{}.

\begin{quote}
\ml{പണ്ട്} \ml{ഗ്രാമങ്ങളിൽ} \ml{കൂട്ടുകുടുംബങ്ങൾ} \ml{സാധാരണമായിരുന്നു}.
\end{quote}

\begin{quote}
\emph{paṇṭŭ grāmaṅṅaḷil kūṭṭukuṭumbaṅṅaḷ sādhāraṇamāyirunnu.}
\end{quote}

\begin{quote}
Long ago joint families were common in the villages.
\end{quote}

\ml{കൂട്ട്} together + \ml{കുടുംബം}: several families under one roof.

\begin{cousinweb}
\textbf{\ml{കൂട്ടുകുടുംബം}} (\emph{kūṭṭukuṭumbaṁ}) is a Malayalam compound built from \ml{കൂട്ട്}, \textquotedblleft{}joining, company\textquotedblright{}, which is native (DEDR 1882), and \ml{കുടുംബം}, from Sanskrit \emph{kuṭumba}, \textquotedblleft{}a household\textquotedblright{} (CDIAL 3233).
\end{cousinweb}

\begin{grammarlens}[title={What you've built}]
One more word for describing what you lived and felt.
\end{grammarlens}

\subsection*{Your turn}
\begin{itemize}
\raggedright
  \item \emph{Say it:} \emph{kūṭṭukuṭumbaṁ}
  \item \emph{Say it:} \emph{kūṭṭukuṭumbaṁ}, once more
  \item \emph{Say it:} \emph{paṇṭŭ grāmaṅṅaḷil kūṭṭukuṭumbaṅṅaḷ sādhāraṇamāyirunnu.} — Long ago joint families were common in the villages.
  \item \emph{Recall:} say the Malayalam for a family member, then the Malayalam for a lullaby, then say \emph{kūṭṭukuṭumbaṁ} again
\end{itemize}

\begin{tcolorbox}[breakable,colback=teal!4,colframe=teal!35!black,title={Before you move on}]
What does \ml{കൂട്ടുകുടുംബം} mean? (\textquotedblleft{}A joint family\textquotedblright{}.) Say it once more.
\end{tcolorbox}
\section[\texorpdfstring{\ml{അണുകുടുംബം} (aṇukuṭumbaṁ)}{അണുകുടുംബം (aṇukuṭumbaṁ)}]{\ml{അണുകുടുംബം} (aṇukuṭumbaṁ) — a nuclear family}

\label{lesson:ML-C555-anukutumbam}



\begin{quote}
Before the new one: say the Malayalam for a lullaby, then the Malayalam for a joint family.
\end{quote}

\subsection*{You'll want to know: \ml{അണുകുടുംബം}}
\textbf{\ml{അണുകുടുംബം}} — \emph{aṇukuṭumbaṁ} — \textquotedblleft{}a nuclear family\textquotedblright{}.

\begin{quote}
\ml{ഇക്കാലത്ത്} \ml{മിക്ക} \ml{ആളുകളും} \ml{അണുകുടുംബമായി} \ml{താമസിക്കുന്നു}.
\end{quote}

\begin{quote}
\emph{ikkālattŭ mikka āḷukaḷuṁ aṇukuṭumbamāyi tāmasikkunnu.}
\end{quote}

\begin{quote}
Nowadays most people live as nuclear families.
\end{quote}

\ml{അണു} atom, something small + \ml{കുടുംബം}: parents and children only.

\begin{cousinweb}
\textbf{\ml{അണുകുടുംബം}} (\emph{aṇukuṭumbaṁ}) is a learned compound built from Sanskrit \emph{aṇu}, \textquotedblleft{}fine, minute\textquotedblright{} (CDIAL 193), and \emph{kuṭumba}, \textquotedblleft{}a household\textquotedblright{} (CDIAL 3233).
\end{cousinweb}

\begin{grammarlens}[title={What you've built}]
One more word for describing what you lived and felt.
\end{grammarlens}

\subsection*{Your turn}
\begin{itemize}
\raggedright
  \item \emph{Say it:} \emph{aṇukuṭumbaṁ}
  \item \emph{Say it:} \emph{aṇukuṭumbaṁ}, once more
  \item \emph{Say it:} \emph{ikkālattŭ mikka āḷukaḷuṁ aṇukuṭumbamāyi tāmasikkunnu.} — Nowadays most people live as nuclear families.
  \item \emph{Recall:} say the Malayalam for a lullaby, then the Malayalam for a joint family, then say \emph{aṇukuṭumbaṁ} again
\end{itemize}

\begin{tcolorbox}[breakable,colback=teal!4,colframe=teal!35!black,title={Before you move on}]
What does \ml{അണുകുടുംബം} mean? (\textquotedblleft{}A nuclear family\textquotedblright{}.) Say it once more.
\end{tcolorbox}
\section[\texorpdfstring{\ml{വേർപിരിയുക} (vērpiriyuka)}{വേർപിരിയുക (vērpiriyuka)}]{\ml{വേർപിരിയുക} (vērpiriyuka) — to split up; to separate}

\label{lesson:ML-C555-verpiriyuka}



\begin{quote}
Before the new one: say the Malayalam for a joint family, then the Malayalam for a nuclear family.
\end{quote}

\subsection*{You'll want to know: \ml{വേർപിരിയുക}}
\textbf{\ml{വേർപിരിയുക}} — \emph{vērpiriyuka} — \textquotedblleft{}to split up; to separate\textquotedblright{}.

\begin{quote}
\ml{പത്ത്} \ml{വർഷത്തിന്} \ml{ശേഷം} \ml{ആ} \ml{ദമ്പതികൾ} \ml{വേർപിരിഞ്ഞു}.
\end{quote}

\begin{quote}
\emph{pattŭ varṣattinŭ śēṣaṁ ā dampatikaḷ vērpiriññu.}
\end{quote}

\begin{quote}
That couple separated after ten years.
\end{quote}

Past: \ml{വേർപിരിഞ്ഞു}, \emph{vērpiriññu}.

\ml{വേർ} apart + \ml{പിരിയുക}, to part, already familiar.

\begin{cousinweb}
\textbf{\ml{വേർപിരിയുക}} (\emph{vērpiriyuka}) is native Dravidian, a compound of \ml{വേറ്}, \textquotedblleft{}separation, different\textquotedblright{} (DEDR 5548), and \ml{പിരിയുക}, \textquotedblleft{}to part\textquotedblright{} (DEDR 4176).
\end{cousinweb}

\begin{grammarlens}[title={What you've built}]
One more word for describing what you lived and felt.
\end{grammarlens}

\subsection*{Your turn}
\begin{itemize}
\raggedright
  \item \emph{Say it:} \emph{vērpiriyuka}
  \item \emph{Say it:} \emph{vērpiriyuka}, once more
  \item \emph{Say it:} \emph{pattŭ varṣattinŭ śēṣaṁ ā dampatikaḷ vērpiriññu.} — That couple separated after ten years.
  \item \emph{Recall:} say the Malayalam for a joint family, then the Malayalam for a nuclear family, then say \emph{vērpiriyuka} again
\end{itemize}

\begin{tcolorbox}[breakable,colback=teal!4,colframe=teal!35!black,title={Before you move on}]
What does \ml{വേർപിരിയുക} mean? (\textquotedblleft{}To split up; to separate\textquotedblright{}.) Say it once more.
\end{tcolorbox}
\section[\texorpdfstring{\ml{ഇണങ്ങുക} (iṇaṅṅuka)}{ഇണങ്ങുക (iṇaṅṅuka)}]{\ml{ഇണങ്ങുക} (iṇaṅṅuka) — to be reconciled; to make up}

\label{lesson:ML-C555-inannuka}



\begin{quote}
Before the new one: say the Malayalam for a nuclear family, then the Malayalam for to split up.
\end{quote}

\subsection*{You'll want to know: \ml{ഇണങ്ങുക}}
\textbf{\ml{ഇണങ്ങുക}} — \emph{iṇaṅṅuka} — \textquotedblleft{}to be reconciled; to make up\textquotedblright{}.

\begin{quote}
\ml{വഴക്കിന്} \ml{ശേഷം} \ml{കുട്ടികൾ} \ml{വേഗം} \ml{ഇണങ്ങി}.
\end{quote}

\begin{quote}
\emph{vaḻakkinŭ śēṣaṁ kuṭṭikaḷ vēgaṁ iṇaṅṅi.}
\end{quote}

\begin{quote}
After the quarrel the children soon made up.
\end{quote}

Past: \ml{ഇണങ്ങി}, \emph{iṇaṅṅi}.

The opposite of \ml{പിണങ്ങുക}, to sulk. It also means to get used to a place or a person.

\begin{cousinweb}
\textbf{\ml{ഇണങ്ങുക}} (\emph{iṇaṅṅuka}) is native Dravidian: DEDR 457 lists it with \ml{ഇണ}, \textquotedblleft{}a pair, a mate\textquotedblright{}, and verbs for suiting and uniting (DEDR 457).
\end{cousinweb}

\begin{grammarlens}[title={What you've built}]
One more word for describing what you lived and felt.
\end{grammarlens}

\subsection*{Your turn}
\begin{itemize}
\raggedright
  \item \emph{Say it:} \emph{iṇaṅṅuka}
  \item \emph{Say it:} \emph{iṇaṅṅuka}, once more
  \item \emph{Say it:} \emph{vaḻakkinŭ śēṣaṁ kuṭṭikaḷ vēgaṁ iṇaṅṅi.} — After the quarrel the children soon made up.
  \item \emph{Recall:} say the Malayalam for a nuclear family, then the Malayalam for to split up, then say \emph{iṇaṅṅuka} again
\end{itemize}

\begin{tcolorbox}[breakable,colback=teal!4,colframe=teal!35!black,title={Before you move on}]
What does \ml{ഇണങ്ങുക} mean? (\textquotedblleft{}To be reconciled; to make up\textquotedblright{}.) Say it once more.
\end{tcolorbox}
